% GDZ PPN236018647,native546–547upper;HKUST547–548.
\clearpage\SourceLeaf{546}\ArithmeticHead{517}{R}{1797 MART.—1797 MART. 29.}
\NumberedItem{62}
\indent Lemniscata geometrice in quinque partes dividitur.
\par\noindent\hfill[1797] M[a]rt. 21.

\medskip\begingroup\small
\indent »Geometrice« heißt mit Zirkel und Lineal. Die Gleichung für die Fünfteilung siehe oben S. 161 ff. artt. [4.] und [5.], die algebraische Darstellung ihrer reellen Wurzeln ebenda und S. 163, art. [6.]. Vergl. auch die bei der Nr. 60 angeführte Stelle der \textit{Disquisitiones arithmeticae}, wobei jedoch zu bemerken ist, daß schon der Graf \textsc{Fagnano} im II. Bande der \textit{Produzioni matematiche} (Pesaro, 1750), S. 356 ff.\footnote{Opere matematiche del \textsc{Marchese G. C. de’ Toschi di Fagnano}, II, 1911, S. 304 ff. \textit{Metodo per misurare la lemniscata}, Schediasma II.} gezeigt hat, daß die Lemniskate sich »algebraisch« in gleiche Teile zerlegen lasse, wenn die Anzahl der Teile in einer der drei Formen \(2.2^m,\ 3.2^m,\ 5.2^m\) enthalten ist, wo \(m\) eine positive ganze Zahl bedeutet. Anknüpfend an die bei der Nr. 60 erwähnte Stelle der \textit{Disquisitiones arithmeticae} und durch Anwendung der von \textsc{Gauss} für die Kreisteilung geschaffenen Methode hat N. H. \textsc{Abel} in seinen \textit{Recherches sur les fonctions elliptiques} (\textsc{Crelles} Journal für Mathematik \textbf{2}, 1827, und \textbf{3}, 1828, Oeuvres, nouvelle édition, I, 1881, S. 261 ff.) nachgewiesen, daß die Teilung der Lemniskate in \(n\) gleiche Teile für diejenigen Werte von \(n\) mit Zirkel und Lineal bewirkt werden kann, für die dies beim Kreise möglich ist (siehe Oeuvres I, S. 314 und 361).
\par\noindent\hfill\textsc{Klein.\quad Schlesinger.}
\par\endgroup

\NumberedItem{63}
\indent Inter multa alia Curvam Lemniscatam spectantia observavi:
\par\indent [1)] Numeratorem sinus decompositi, arcus duplicis esse \(=\)
\par\noindent 2 Num. Denom. Sinus \(\times\) Num. Den[om]. Cos. arcus simpl[icis];
\par\indent [2)] Denominatorem vero \(=(\text{Num. sin.})^4+(\text{Denom. sin.})^4\).
\par\indent [3)] Iam si hic denominator pro arcu \(\pi^l\) ponatur \(\theta\), erit Denom[inator] sin arcus \(k\pi^l,=\theta^{kk}\).
\par\indent [4)] Iam
\[
\theta=4{,}810480,
\]
\indent [5)] cuius numeri logarithmus hyperbolicus est
\[
=1{,}570796\quad\text{i. e.}\ =\tfrac12\pi,
\]
\noindent quod maxime est memorabile, cuiusque proprietatis demonstratio gravissima analyseos incrementa pollicetur.
\par\noindent\hfill[1797] Mart. 29.

\medskip\begingroup\small
\indent Man findet die hier zusammengestellten fünf Aussagen in den oben S. 156 ff. artt. [7.] und [8.] abgedruckten Leisteaufzeichnungen entwickelt. Setzt man wie a. a. O.
\[
\operatorname{sin\,lemn}\varphi=\frac{M(\varphi)}{N(\varphi)},\quad\operatorname{cos\,lemn}\varphi=\frac{\mu(\varphi)}{\nu(\varphi)},
\]
\par\endgroup

\clearpage\SourceLeaf{547}\ArithmeticHead{518}{L}{TAGEBUCH MIT ERLÄUTERUNGEN.}
\begingroup\small
\noindent so ist nach den beiden ersten Aussagen (siehe oben S. 157)
\[
\begin{aligned}
M(2\varphi)&=2M(\varphi)N(\varphi)\mu(\varphi)\nu(\varphi),\\
N(2\varphi)&=M(\varphi)^4+N(\varphi)^4.
\end{aligned}
\]
\indent Die hier in der Tagebuchaufzeichnung mit \(\pi^l\) bezeichnete Größe wird in der Leistehandschrift mit \(\Pi\) bezeichnet, wir haben dafür in dem Abdruck das von \textsc{Gauss} später stets benutzte \(\varpi\) gesetzt. Die hier mit \(\theta\) bezeichnete Größe heißt im \textsc{Leiste} (siehe oben S. 157) \(2a^4\); es ist also nach der dritten Aussage (siehe oben S. 157)
\[
N(\varpi)=\theta,\quad N(k\varpi)=\theta^{kk}
\]
\noindent für ein ganzes positives \(k\). Ferner ist nach der vierten und fünften Aussage (siehe oben S. 158)
\[
N(\varpi)=4{,}81048,\quad\log4{,}81048=1{,}5708=\frac\pi2.
\]
\noindent Die »demonstratio« der hier nur durch numerische Rechnung erkannten Eigenschaft [5)] gelingt \textsc{Gauss} erst im Juli 1798; vergl. die Bemerkungen zu der Tagebuchaufzeichnung Nr. 92.
\par\noindent\hfill\textsc{Klein.\quad Schlesinger.}
\par\endgroup
